\section{Studies}
\label{Sec:Studies}
We now illustrate our sequential Bayesian meta-analysis framework. Our
goal is to illustrate how the method tracks the evolution of
scientific evidence, quantifies the impact of new studies, and
explicitly models biases. Specifically, we address these
methodological questions: How do beliefs about a scientific parameter
evolve as studies accumulate? How does explicit modeling of systematic
biases influence our conclusions? How sensitive are results to our
prior assumptions?

We will begin with a large-scale replication dataset from psychology (Many Labs, Section \ref{Sec:ManyLabs}). SMART recovers the diminishing returns on research that are well-known for genuine (and successful) replications. 

We then move to a simulated example (Section \ref{Sec:ToyExample}) that illustrates the impact of methodological advances on a literature's trajectory. The advent of a new method may lead to a reevaluation of beliefs about past studies' biases.
 SMART picks up on the special role of pioneering studies and shows how their influence depends on a research community's beliefs (and their changes) about the biases of its methods.

We end with a case study in which we explore the importance of Card and Krueger's seminal 1994 study \citep{Card_Krueger_94}, which, due to its persuasive design, was essential in challenging prevailing economic thought on the effects of a minimum wage on employment.

\subsection{Many Labs}
\label{Sec:ManyLabs}
\begin{figure*}[t!]
    \includegraphics[width=\textwidth]{figure/ManyLabs.pdf}
   
   \caption{Sequential meta-analysis trace for the Many Labs example. (Left) Sequence of study outcomes replicating the ``Imagined Contact'' study \citep{Klein_2014}. (Middle) Evolution of the posterior mean with its 95\% credible interval. (Right) Sequential research contribution of every study, measured by the Wasserstein-2 learning metric. The trace recovers the diminishing returns that are to be expected when a single method is applied repeatedly.}
  \label{fig:ML_W2}
\end{figure*}
We begin by analyzing the Many Labs
dataset, a large-scale psychology replication project, which allows us
to track how evidence about certain psychological effects evolves as
studies from multiple laboratories accumulate. We analyze the original
Many Labs data, illustrating sequential posterior updating in
a real replication scenario. 
Data and code for all
examples may be found on \href{https://github.com/JonasMikhaeil/SequentialMATrace}{github}.

Each Many Labs study contains a collection of survey responses. The
design is described in \citep{Klein_2014}, and we focus on the
``Imagined Contact'' study. In the treatment condition, subjects are
asked to imagine interacting with a Muslim stranger for one
minute. (One exception is when the study is in Turkey, where they are
asked to imagine interacting with a Christian stranger.) In the control
condition, subjects are asked to imagine that they are walking
outdoors for the same length of time. After imagining their assigned
scene, subjects respond to a series of questions. The outcome variable
is a composite index of four questions that ask about interest in
learning about Islam and interacting with Muslims.

Specifically, each response is on a 1-9 scale, with higher values
indicating a more positive attitude, and the outcome is the average of
these four items. The studies were designed to be as similar as
possible in terms of the stimulus and outcome measures. So they
comport with the RE meta-analysis model, which attributes
study-to-study variation as a combination of statistical error and
idiosyncratic effect heterogeneity. Our goal is to estimate the
expected difference in outcome between treatment and control.

Figure \ref{fig:ML_W2} shows the sequential research trace for the
original sequence of studies. The initial studies contribute
substantially to our shared knowledge and significantly affect the
collective posterior. Eventually, however, the marginal contribution
of each study tapers off, and studies later in the sequence have
little effect on the posterior.

\subsection{A Simulation of Methodological Innovation}
\label{Sec:ToyExample}


Advances in research methodology can lead to surprising results. In
this section, we present a simulated example that illustrates how our
method tracks nonlinear scientific progress.  We simulate data in which
study outcomes $y_i$ are estimates of a parameter of interest
$\vartheta$, but contaminated by a method-specific bias $z_i$.

Specifically, we draw data from the following data-generating process
\begin{align}
   z_i \lvert \ell_i \, &\sim  \, \text{normal}(\beta \ell_i,.01) \\ \nonumber
   y_i \lvert z_i \, &\sim \, \text{normal}(\vartheta^* + z_i,.01).    
\end{align}
Here correlations in the research methodology of different labs
$\ell_i$ lead to correlations of the results $y_i$. We simulate a
scenario in which the first ten experiments ($1 \leq i \leq 10$) use a
specific and established methodology $\ell_i = 1$. After a while,
researchers develop a new and more accurate methodology $\ell_i = 0$,
and use it for the subsequent experiments ($11 \leq i \leq 30$). Figure \ref{fig:syn_trace} (Left) shows a sequence of 
experimental outcomes for this example.

Using the methodological labels, we analyze this data
with the labeled random effects model of
Section \ref{Sec:QuantifyingLearning}\ref{sec:random-effects-model}. What role do our beliefs
about the methods' accuracy (before and after the change of
methodology) play?

To demonstrate these methods, we consider two possible scientific
scenarios.  In the first scenario (I), we suspect the original
methodology to be biased (say, $\kappa_1 = 1$). Through conceptual
advances, the pioneering study (No. 11) introduces a new methodology
that the research community holds to be unbiased
($\kappa_2 = 0.0001 \ll \theta$). Figure \ref{fig:syn_trace}
shows the sequential meta-analysis trace for this example.
The middle panel tracks the evolution of the posterior. The right panel shows how much is learned by each study measured in terms of our
learning metric. The methodological shift to an unbiased method makes
the 11\textsuperscript{th} study highly influential.

In a second scenario (II), suppose that the first methodology is
originally believed to be accurate ($\kappa_{i<11,1} =
0.0001$). Theoretical and conceptual research, however, uncovers flaws
in this methodology, which eventually leads the research community to
acknowledge the possibility of bias for this methodology
($\kappa_{t\geq 11,1} = \kappa_1$). Building on this development, a
pioneering study (No. 11) introduces a new method into the literature,
and this method is believed to yield unbiased results
($\kappa_{i\geq 11,2} = 0.0001$).  Figure \ref{fig:sens_syn} tracks the posterior for different degrees $\kappa_1$ to which the research community
begins to doubt the first methodology. The middle and right panel show that posterior uncertainty can
increase when a new research methodology establishes itself even if
beliefs about methodological biases are not taken into
account. Studies that pioneer a change in methodology can be
pathbreaking, because they cast doubt on what has previously been
learned. The figure also highlights how much is learned from the pioneering study (No. 11)
depends on the degree $\kappa_1$ to which the research community
begins to doubt the first methodology.

Both scenarios highlight the importance of innovation in
research. Through the development of more accurate methods, there can
be a high return to a study in terms of learning even when the same
literature would face diminishing returns sticking to a single
method. Our second scenario highlights the contribution of theoretical
work (that may not contribute any data at all) to the importance of
pioneering empirical studies by informing our beliefs about a
literature's potential biases

Finally, we use this simulation to illustrate the difference between
our sequential meta-analysis trace and classical meta-analysis
weights \citep{Riley_Ensor_Jackson_Burke_2018,Burke_2018}.\footnote{For a definition of meta-analysis weights, and a more in-depth discussion, see Appendix \ref{Sec:CompMAW_SMAT}.} Table \ref{tab:MAweights} (in Appendix \ref{Sec:CompMAW_SMAT}) contains
meta-analysis weights calculated sequentially and retrospectively for our stylized example. Despite pioneering novel methodology, study 11 receives no outstanding weight.
Through the lens of meta-analysis, the value of a study reflects its contribution to the overall meta-analysis effect estimate. 
No account is taken of whether certain studies may have been ground-breaking by advancing new methodologies. 
When meta-analysis weights are used to gauge a study's contribution to the literature, studies with small reported standard errors receive higher weights. This essentially favors large (and oftentimes expensive) studies while diminishing the importance of smaller and innovative pilot studies. Through the lens of our method for quantifying learning, the precision of a study is only one factor contributing to its impact. By being able to quantify the value of innovation, our method gives due credit to pioneers. 
\begin{figure*}[t!]
    \includegraphics[width=\textwidth]{figure/SynExmp.pdf}
  \caption{Sequential meta-analysis trace for the stylized example of methodological innovation. (Left) Sequence of experimental outcomes. The first 10 studies are biased. The 11th study pioneers improved and unbiased methodology that takes hold in the literature. (Middle) Evolution of the posterior mean and its uncertainty ($95$\% credible interval). Because the original methodology is believed to be biased, the posterior remains diffuse for the first 10 studies. By introducing an unbiased methodology, the 11th study shifts the posterior towards the true effect and significantly reduces posterior uncertainty. (Right) Sequential research contribution measured with the Wasserstein-2 learning metric. Due to its impact on the posterior both in terms of its location and scale, the 11th is highly influential. Our trace picks up on its significant research contribution.}
  \label{fig:syn_trace}
\end{figure*}
\begin{figure*}[t!]
    \includegraphics[width=\textwidth]{figure/Sensitivity_ToyExample.pdf}
    \caption{Evolution of the posterior for different beliefs about methodological biases. In all three situations, the methodology of the first 10 studies is originally believed to be unbiased ($\kappa_{i<11,1} =
0.0001$). The 11th study pioneers new unbiased methodology and advances theoretical reasons that cast doubt on the original methodology, changing the community's beliefs about the method ($\kappa_{i\geq 11,1} = \kappa_1$). The panels show the evolution of the posterior in the case that the beliefs change to (Left) $\kappa_1=1$, (Middle) $\kappa_1=0.2$ or (Right) remain unchanged $\kappa_1=0.0001$.}
  \label{fig:sens_syn}
\end{figure*}

\subsection{Case Study: Assessing the Employment Effects of Minimum Wage Laws}
\label{Sec:CaseStudy}

\begin{figure*}[t!]
    \includegraphics[width=\textwidth]{figure/MW_plot.pdf}
    \caption{Sequential meta-analysis trace for the minimum wage literature leading up to Card and Krueger's 1994 study \citep{Card_Krueger_94}. 
    We begin with a prior $\pi_0(\theta) = \mathcal{N}(-0.15,0.07^2+0.3^2)$ based on discounted time-series evidence \citep{Brown_82}. We assume that the methodologies used before C\&K's study were believed to be biased ($\kappa_{1} = \kappa_{2} = \kappa_{3} = 0.3$) and that C\&K's border discontinuity design was believed to be less biased ($\kappa_4 = 0.05)$. 
    (Left) Sequence of estimates leading up to C\&K's study (with their estimate in
      orange). Estimates colored in blue belong to the same study \citep{Williams_93} and the posterior is updated based on them collectively. (Middle) Evolution of the posterior mean and its $95$\% credible interval. (Right) Sequential research contribution measured with the Wasserstein-$2$ learning metric. This trace recovers the influential role C\&K's study played in the New Minimum Wage Research \citep{Ehrenberg_92} literature.  }
  \label{fig:MW_figure}
\end{figure*}
\begin{figure}[th]
\centering
\includegraphics[width=0.7\textwidth]{figure/sensanalysis_MW.pdf}
\caption{Sensitivity analysis for C\&K's 1994 study's \citep{Card_Krueger_94} sequential research contribution. Their contribution crucially depends on how unconfounded their border discontinuity design is believed to be. For values of $\kappa_4$ close to $0$ C\&K's results is credible. With increases in credibility, a more significant sequential research contribution is accorded to C\&K's study.}
\label{fig:MW_sensitivity}
\end{figure}
In this section, we present a real-world illustration of our method to
assess sequential research contributions. We focus on a literature estimating the effect of a minimum wage on employment and, specifically, seek to understand the role of the seminal 1994 study by Card
\& Krueger \citep{Card_Krueger_94} --- we will refer to them as C\&K --- who
used the 1992 increase in New Jersey's minimum wage to estimate
this effect. 
The minimum wage is an important policy lever intended to increase income among low-wage
earners. Economists worry, however, that raising the minimum wage may negatively affect
employment, hurting some low-wage workers while helping others. 
The best evidence prior to the 1990s, which was based on time-series data, suggested that increases in the minimum wage led to a reduction in employment, a result in line with the textbook microeconomics.
Labor economists by the late 1980s, however, increasingly discounted these empirical results because of potential confounding.
The 1990s saw the ascent of a \textit{New Minimum Wage Research} literature
\citep{Ehrenberg_92}, featuring innovative study designs that would challenge classical economic thought on this topic \citep{myth_95}. 

We look at the beginning of this literature, which featured a methodological shift from
time-series analysis to estimation strategies leveraging plausibly exogenous variation in the minimum wage. 
The goal of the new
approach was to reduce the threat of omitted variables bias and thus provide more credible estimates. Specifically, we analyze the trajectory of six
studies from 1992 to 1994 leading up to, and including, C\&K as presented by Neumark and Shirley
\citep{Neumark_Shirley_22}.\footnote{
All estimates report elasticities. For studies reporting multiple estimates of the same effect with different controls \citep{Cardb_1992,Neumark_1992}, we selected one representative estimate.} Figure \ref{fig:MW_figure} shows the sequential research trace for this literature. 

One of the early contributions to this new research literature was C\&K's
classic study \citep{Card_Krueger_94}, which rose to prominence because of its convincing design. To mitigate confounding, the authors applied 
difference-in-differences (DID) estimation to employment data gathered along a border discontinuity, comparing counties in New Jersey affected by the increase of the state's minimum wage with adjacent unaffected counties in Pennsylvania  \citep{nobel_21}. Another compelling aspect of this design was its prospective collection of data, which reduced researcher degrees of freedom.\footnote{Card was awarded
  the Nobel Memorial Prize in Economic Sciences partly for this and
  related work. DID itself was introduced to the minimum
  wage literature in 1915 by Obenauer and von
  der Nienburg \citep{Obenauer_Nienburg}, see
  \citep{Angrist_Pischke}. Card also used DID to estimate the effect of a minimum wage increase in California \citep{Card_92}.}

There is widespread agreement that C\&K's study was highly influential.\footnote{To give just one quotation
  highlighting C\&K's study perceived influence: ``The path-breaking
  work of Card and Krueger [1994], showing that a higher minimum wage
  can increase employment, turned the age-old conventional wisdom on
  its head.''  \citep{marjit_21}} Their study has become a widely-cited
textbook example of a well-designed study with enormous policy impact \citep{Angrist_Pischke}. It also inspires ongoing
theoretical \citep{Dube_16,Card_18,marjit_21,Renkin_22} and methodological advances
\citep{Neumark_13,Abadie_2005,Athey_2006,torous2024optimaltransportapproachestimating}.

Classical meta-analysis, however, accords it little weight. Applying the random-effects model to the sequence of studies from 1992 to 1994, we find that C\&K's study is given a weight of 15\%, which is the second smallest weight in this collection of studies. Unlike classical meta-analysis, the SMART analysis in Figure \ref{fig:MW_figure} picks up on the seminal impact of C\&K.


Our aim is not to weigh in on the employment effects of the minimum
wage, which continue to be the focus of research and
debate \citep{Neumark_Shirley_22,Manning_21}. Rather, we use our
sequential meta-analysis trace to understand the influence of C\&K's study on the development of
this literature. 

What made C\&K's study so influential?  First, C\&K's study found a surprising result. This alone,
however, would not have been enough. It is only in conjunction with
the research community's widespread belief that
C\&K's methodology is more accurate than previous methods that
their surprising results shifted the collective posterior.

To formulate a research trace for this literature, we begin by reconstructing a prior. The time-series evidence up to the early 1980s is summarized in \citep{Brown_82}, which concludes that negative elasticities between 1$\%$ and 3$\%$ are most plausible with the lower range being more likely. A rough approximation of this summary is a prior such as $\mathcal{N}(-0.15,0.07^2)$.  
A growing concern with confounding (see, for example, \citep{LaLonde_86})  led labor economists in the late 1980s to discount this time-series evidence. Somewhat arbitrarily, we can widen the prior $\pi_0(\theta) = \mathcal{N}(-0.15,0.07^2+0.3^2)$ to express this additional uncertainty, but note that the story our trace conveys does not depend on the exact illustrative numbers chosen.

To quantify the influence of C\&K's study, we again use the labeled
random-effects model of Section \ref{Sec:QuantifyingLearning}\ref{sec:random-effects-model}.  We group all methodologies using panel data on federal variation of the minimum wage under $\ell=\ell_{fed}$ \citep{Cardb_1992,Katz_92,Williams_93} and those using state-level variation under $\ell=\ell_{st}$ \citep{Neumark_1992}. $\ell=\ell_{DID}$ denotes the DID methodology used in \citep{Card_92} and $\ell =\ell_{CK}$ indicates C\&K's border discontinuity design \citep{Card_Krueger_94}.

While the studies based on (observational) panel data were all aiming to leverage plausibly exogenous variation in the minimum wage, the threat of omitted variables bias remained a lingering concern \citep{nobel_21}.
For simplicity, we assume that all studies predating C\&K were believed to be biased to the same degree $\kappa_{1} = \kappa_{2} = \kappa_{3} = 0.3$. To determine the effect of the minimum wage on employment, the effects of the change in the minimum wage need to be disentangled from other factors affecting employment. 
By comparing adjacent counties that experience similar economic shocks, some affected by a rise in the minimum wage while others remained unaffected, C\&K's border discontinuity design 
plausibly mitigated confounding. This innovative design led to the belief that this study was less biased than its predecessors, say, $\kappa_4= 0.05$.

Figure \ref{fig:MW_figure} showcases the significance of C\&K's study. 
On the left, we see the sequence of estimates leading up to C\&K's study. The middle panel tracks how the posterior is sequentially updated with every study. The plot on the right shows each study's sequential research contribution as it enters the literature, with C\&K's study clearly standing out for its importance.

Figure \ref{fig:MW_sensitivity} provides a sensitivity analysis for assessing the importance of C\&K's study
\citep{Card_Krueger_94}. It shows how C\&K's study's research contribution decreases as a function of $\kappa_{4}$: the more convincingly their design overcomes confounding, the more influential their study becomes. If \citep{Card_Krueger_94} were just another flawed study, its influence would have been unremarkable.





